% GENERATED FILE. Edit canonical lessons, then run npm run generate:books.
% canonical-source-hash: fnv1a64:e4d5fd74fe23593c
% canonical-lessons: PA-C133-kukkar, PA-C133-murgi, PA-C133-gadha, PA-C133-ullu, PA-C133-ka

\chapter{Rooster, Hen, Donkey, Owl, Crow}
\label{ch:pa-rooster-hen-donkey-owl-crow}
\hlchaptermodality{133}

\begin{chapteropening}
\textbf{By the end of this chapter:} \emph{I can say a rooster, a hen, a donkey, an owl and a crow.}

Complete the last lesson of chapter 133: I can say a rooster, a hen, a donkey, an owl and a crow.
\end{chapteropening}

\section[\texorpdfstring{\pa{ਕੁੱਕੜ} (kukkaṛ)}{ਕੁੱਕੜ (kukkaṛ)}]{\pa{ਕੁੱਕੜ} (kukkaṛ) — a rooster}

\label{lesson:PA-C133-kukkar}



\begin{quote}
Before the new one: say the Punjabi for a spider, then the Punjabi for a squirrel.
\end{quote}

\subsection*{You'll want to know: \pa{ਕੁੱਕੜ}}
\textbf{\pa{ਕੁੱਕੜ}} — \emph{kukkaṛ} — \textquotedblleft{}a rooster\textquotedblright{}.

\begin{grammarlens}[title={What you've built}]
One more word for this chapter.
\end{grammarlens}

\subsection*{Your turn}
\begin{itemize}
\raggedright
  \item \emph{Say it:} \emph{kukkaṛ}
  \item \emph{Say it:} \emph{kukkaṛ}, once more
  \item \emph{Say it:} say \emph{kukkaṛ}
  \item \emph{Recall:} say the Punjabi for a spider, then the Punjabi for a squirrel, then say \emph{kukkaṛ} again
\end{itemize}

\begin{tcolorbox}[breakable,colback=teal!4,colframe=teal!35!black,title={Before you move on}]
What does \pa{ਕੁੱਕੜ} mean? (\textquotedblleft{}A rooster\textquotedblright{}.) Say it once more.
\end{tcolorbox}
\section[\texorpdfstring{\pa{ਮੁਰਗੀ} (murgī)}{ਮੁਰਗੀ (murgī)}]{\pa{ਮੁਰਗੀ} (murgī) — a hen}

\label{lesson:PA-C133-murgi}



\begin{quote}
Before the new one: say the Punjabi for a squirrel, then the Punjabi for a rooster.

\emph{Read it:} \textbf{\pa{ਕਿੱਥੇ}}

What does the small mark before \textbf{\pa{ਥ}} do? (The addak doubles it: \emph{kit-the}, \textquotedblleft{}where.\textquotedblright{})
\end{quote}

\subsection*{You'll want to know: \pa{ਮੁਰਗੀ}}
\textbf{\pa{ਮੁਰਗੀ}} — \emph{murgī} — \textquotedblleft{}a hen\textquotedblright{}.

\begin{grammarlens}[title={What you've built}]
One more word for this chapter.
\end{grammarlens}

\subsection*{Your turn}
\begin{itemize}
\raggedright
  \item \emph{Say it:} \emph{murgī}
  \item \emph{Say it:} \emph{murgī}, once more
  \item \emph{Say it:} say \emph{murgī}
  \item \emph{Recall:} say the Punjabi for a squirrel, then the Punjabi for a rooster, then say \emph{murgī} again
\end{itemize}

\begin{tcolorbox}[breakable,colback=teal!4,colframe=teal!35!black,title={Before you move on}]
What does \pa{ਮੁਰਗੀ} mean? (\textquotedblleft{}A hen\textquotedblright{}.) Say it once more.
\end{tcolorbox}
\section[\texorpdfstring{\pa{ਗਧਾ} (gadhā)}{ਗਧਾ (gadhā)}]{\pa{ਗਧਾ} (gadhā) — a donkey}

\label{lesson:PA-C133-gadha}



\begin{quote}
Before the new one: say the Punjabi for a rooster, then the Punjabi for a hen.
\end{quote}

\subsection*{You'll want to know: \pa{ਗਧਾ}}
\textbf{\pa{ਗਧਾ}} — \emph{gadhā} — \textquotedblleft{}a donkey\textquotedblright{}.

\begin{grammarlens}[title={What you've built}]
One more word for this chapter.
\end{grammarlens}

\subsection*{Your turn}
\begin{itemize}
\raggedright
  \item \emph{Say it:} \emph{gadhā}
  \item \emph{Say it:} \emph{gadhā}, once more
  \item \emph{Say it:} say \emph{gadhā}
  \item \emph{Recall:} say the Punjabi for a rooster, then the Punjabi for a hen, then say \emph{gadhā} again
\end{itemize}

\begin{tcolorbox}[breakable,colback=teal!4,colframe=teal!35!black,title={Before you move on}]
What does \pa{ਗਧਾ} mean? (\textquotedblleft{}A donkey\textquotedblright{}.) Say it once more.
\end{tcolorbox}
\section[\texorpdfstring{\pa{ਉੱਲੂ} (ullū)}{ਉੱਲੂ (ullū)}]{\pa{ਉੱਲੂ} (ullū) — an owl}

\label{lesson:PA-C133-ullu}



\begin{quote}
Before the new one: say the Punjabi for a hen, then the Punjabi for a donkey.
\end{quote}

\subsection*{You'll want to know: \pa{ਉੱਲੂ}}
\textbf{\pa{ਉੱਲੂ}} — \emph{ullū} — \textquotedblleft{}an owl\textquotedblright{}.

\begin{grammarlens}[title={What you've built}]
One more word for this chapter.
\end{grammarlens}

\subsection*{Your turn}
\begin{itemize}
\raggedright
  \item \emph{Say it:} \emph{ullū}
  \item \emph{Say it:} \emph{ullū}, once more
  \item \emph{Say it:} say \emph{ullū}
  \item \emph{Recall:} say the Punjabi for a hen, then the Punjabi for a donkey, then say \emph{ullū} again
\end{itemize}

\begin{tcolorbox}[breakable,colback=teal!4,colframe=teal!35!black,title={Before you move on}]
What does \pa{ਉੱਲੂ} mean? (\textquotedblleft{}An owl\textquotedblright{}.) Say it once more.
\end{tcolorbox}
\section[\texorpdfstring{\pa{ਕਾਂ} (kā̃)}{ਕਾਂ (kā̃)}]{\pa{ਕਾਂ} (kā̃) — a crow}

\label{lesson:PA-C133-ka}



\begin{quote}
Before the new one: say the Punjabi for a donkey, then the Punjabi for an owl.
\end{quote}

\subsection*{You'll want to know: \pa{ਕਾਂ}}
\textbf{\pa{ਕਾਂ}} — \emph{kā̃} — \textquotedblleft{}a crow\textquotedblright{}.

\begin{grammarlens}[title={What you've built}]
One more word for this chapter.
\end{grammarlens}

\subsection*{Your turn}
\begin{itemize}
\raggedright
  \item \emph{Say it:} \emph{kā̃}
  \item \emph{Say it:} \emph{kā̃}, once more
  \item \emph{Say it:} say \emph{kā̃}
  \item \emph{Recall:} say the Punjabi for a donkey, then the Punjabi for an owl, then say \emph{kā̃} again
\end{itemize}

\begin{tcolorbox}[breakable,colback=teal!4,colframe=teal!35!black,title={Before you move on}]
What does \pa{ਕਾਂ} mean? (\textquotedblleft{}A crow\textquotedblright{}.) Say it once more.
\end{tcolorbox}
